% =====================================================================
%  Thesis assignment (consolidated) — ENGLISH (translation of de.tex)
%  General scope (search algorithms/trees), design-space vocabulary with
%  citation; NO axis list, NO organism/organ metaphor (that is the solution,
%  linear build-up). Sources: docs/termine (Termin 1-9) + Idreos design space.
%  German master: de.tex.
% =====================================================================
{\large\textbf{Thesis Assignment}}\\[0.4em]
{\itshape Active Cache-Aware Hardware Adaptation:\\ A Cache Engine for Trie-Based Index Structures}

\bigskip
\minisec{Context and motivation}
Search algorithms and index structures are central to modern in-memory database systems, and the
Chair of Databases focuses explicitly on hardware-conscious optimisation of such systems. Since the
Adaptive Radix Tree, numerous cache-aware techniques and related algorithm families have been
proposed --- from node layout through prefetching and succinct representations to allocation ---,
yet almost every publication uses its own data structures, different load profiles, and inconsistent
measurement methodology, which makes results hard to compare across publications. As a consequence,
the weightings of the improvements of individual algorithm components are currently not separable
from one another.

\minisec{Problem statement}
Search methods --- from balanced search trees (B-/B\textsuperscript{+}-trees) through tries and
radix trees to specialised, custom-built constructs --- span a large \emph{design space}: even a
two- or three-component data structure admits astronomically many designs, and a concrete algorithm
is a single point in this space, described by many individual design decisions --- how nodes are
represented, traversed, laid out in memory, and allocated
\cite{idreos2018periodic, idreos2018datacalculator}. Cache-line-aware behaviour arises not at a
single one of these places but at \emph{many different structural places} of an algorithm and from
their interaction. The state of the art comprises numerous mature methods, yet each optimises only
individual ones of these places and demonstrates its benefit in its own measurement environment,
which makes progress hard to compare across methods. A common, generic concept and
\textbf{principle} is missing that can measure and evaluate cache-line-aware behaviour across the
different places of a search algorithm --- interchangeably, systematically, and fairly.

\minisec{Objective}
The aim of the thesis is to find, design, and prototype --- as an experimentation framework --- such
a principle (a framework) that (i) decomposes cache-aware search algorithms into interchangeable
design components, (ii) reconstructs state-of-the-art algorithms as explicit configurations (points
in the design space), (iii) measures and evaluates cache-line-aware behaviour both at the individual
directly and indirectly affected places and for the whole algorithm on a common, CPU-only basis, and
(iv) on the basis of the fine-grained measurement results, delivers the best overall algorithm as a
standalone binary behind a unified search-algorithm interface. As a probe for the trial integration
of a new technology into the state of the art, the thesis contributes \mbox{PRT-ART} (a Probst
Redirect Tree~/ ART hybrid), which is studied as a conceptually novel search algorithm, yet
contributes only a few new aspects within the classification of search algorithms and not a complete
algorithm.

The separability of the individual algorithm components demanded here is realized in this thesis by
the axis framework (Chapter~\ref{ch:concept}) and the three-granularity measurement methodology
(Chapter~\ref{ch:methodology}).

\minisec{Sub-tasks}
\begin{enumerate}
  \item Systematically capture the general design components of search algorithms --- in particular
  the cache-aware ones --- and identify the places at which cache-line-aware and
  search-algorithm-specific decisions take effect.
  \item Reconstruct state-of-the-art designs --- where possible from the original code of the
  respective publication --- (ART, HOT, Masstree, CoCo-trie, START, B\textsuperscript{2}-tree,
  Wormhole, SuRF, and further methods) as explicit, comparable configurations.
  \item Design and implement PRT-ART as a probe with several design components of its own, competing
  against the captured state of the art.
  \item Build a reproducible measurement pipeline (binary~$\to$~CSV~$\to$~LaTeX~$\to$~diagram) with
  profile-aware YCSB workload routing over \emph{all} available workloads and hardware-counter
  collection.
  \item Measure systematically and as completely as possible across the configurations in three
  mandatory measurement series --- (A)~probe vs.\ the state of the art, (B)~systematic variation of
  the design decisions, (C)~merge/regression of old vs.\ new variants ---, with each series
  collected at three granularities: \emph{micro-benchmarking} (each individual design component
  within the search-algorithm shell is measured in isolation, through the common interface of its
  category, against load profiles), \emph{macro-benchmarking} (the individual interface operations
  of the search algorithm against sample loads, e.g.\ insert, point lookup, delete), and
  \emph{whole-system benchmarking} (established YCSB load profiles over every recombination of the
  design components).
  \item Evaluate at which places and in which combinations cache-line-aware decisions are beneficial
  or detrimental for which data-type-specific load patterns --- derived from the ranking across the
  three benchmarking granularities and presented as analysable diagrams and tables.
  \item Automation of system detection and binary selection of the most optimal cache-line-aware
  algorithm, based on empirically substantiated measurements, together with a measurement-driven,
  dynamically controlled cache-line-aware adaptation of the algorithm behaviour as a static
  heuristic system policy for the algorithm behaviour under varying loads and operations. This
  so-called extraction is to be stored as an XML load-profile result for the testing environment of
  an overall algorithm, under reference to the various applied load profiles, and is reusable for
  each architecture focus in order to generate heuristics automatically.
\end{enumerate}

\minisec{Method and evaluation}
The evaluation covers, besides the CPU-only core, SIMD and concurrent multi-threading paths, and
compares the configurations against established baselines, including a non-tree-like hash table
(\texttt{flat\_hash\_map}) as a counterprobe within the same genus. The data basis consists of real
string datasets of differing characteristics --- e.g.\ \texttt{url}, \texttt{dna}, \texttt{protein},
\texttt{xml}, \texttt{tpcds-id}, and \texttt{trec-terms} --- under varied data volumes, key lengths,
prefix overlap, local density, skew, hit/miss ratios, and value sizes. Measured quantities include
exact/prefix lookup, prefix enumeration, insert/update/delete, p50/p95/p99 latency, throughput,
bytes per key, and hardware counters (cycles, instructions, L1/L2/LLC and dTLB misses, branch
mispredictions). Single-thread behaviour is studied first, followed by read-parallel multi-thread
scenarios.

Fair comparability is ensured by fixed rules: a common minimal mode (common-denominator) is
separated from the PRT-ART native mode, all methods receive identical key spaces and seeds, missing
capabilities are reported as N/A rather than covertly emulated, and the compiler, flags, ISA path,
allocator, and commit hash of each baseline are logged. For reproducibility, every configuration is
measured in several separate runs whose spread remains part of the statement (percentiles via HDR
histograms rather than averaging); each dataset carries a record of source, checksum, and seed rule.

\minisec{Scope}
The core contribution is CPU-focused with SIMD extensions and multi-threading, but restricted to
read-parallel methods with serialised writing (multi-read, single-write). The GPU is for now out of
scope (extendable in the future); full write parallelism and an engine/optimiser integration are
optional extensions. All architecture extensions are themselves dynamic extensions against the CPU
baseline run and lead to deeper measurement permutations with multiple runtimes. The
system-optimized binary can be built optionally with or without the measurement observer --- as a
debug, measurement, or release variant.
